\chapter{Introduction}
\label{ch:intro}

\section{Why procurement disputes are worth modelling}
\label{sec:context}

UK public bodies spend somewhere around \pounds300 billion a year on goods, works
and services \cite{cabinetoffice2020}, and in 2023 the rulebook for how they spend
it changed. The Procurement Act 2023 \cite{pa2023} swept away a patchwork of
EU-derived regulations and replaced it with a single regime built around four
statutory objectives in section 12: value for money, maximising public benefit,
sharing information for transparency, and acting with integrity. It also
tightened what happens to losing bidders. Section 50 requires a contracting
authority to send each unsuccessful supplier an assessment summary explaining how
its tender was scored; section 51 then opens an eight-working-day standstill
period, starting the day the contract award notice is published, during which the
contract cannot be signed. A supplier who thinks the evaluation went wrong has
that window to raise it, and thirty days from when it knew or ought to have known
of the breach to actually start proceedings in the Technology and Construction
Court.

Those deadlines are what make this problem interesting rather than merely legal.
Figure~\ref{fig:lifecycle} shows the shape of it. Nearly all the leverage sits in
the first fortnight, before anyone has instructed a solicitor, and nearly all the
information asymmetry sits there too. The authority knows exactly what its
evaluation panel did; the losing bidder knows only whatever the assessment
summary was willing to tell it. Once the automatic suspension kicks in and the
case reaches the TCC, both sides are locked into a process that, across this
project's own case corpus, took anywhere from four months to two years to
resolve.

\begin{figure}[htbp]
\centering
\resizebox{\textwidth}{!}{%
\begin{tikzpicture}[font=\footnotesize,>=Latex]
 % timeline spine
 \draw[line width=1.1pt,f21ink] (0,0) -- (13.2,0);
 \foreach \x in {0,2.6,5.2,8.4,11.6} \draw[f21ink,line width=1.1pt] (\x,-0.13) -- (\x,0.13);

 % labels below
 \node[anchor=north,text width=2.3cm,align=center] at (0,-0.25)
 {Evaluation\\ complete};
 \node[anchor=north,text width=2.4cm,align=center] at (2.6,-0.25)
 {Assessment\\ summary sent\\ \textit{(s.50)}};
 \node[anchor=north,text width=2.4cm,align=center] at (5.2,-0.25)
 {Standstill ends\\ 8 working days\\ \textit{(s.51)}};
 \node[anchor=north,text width=2.5cm,align=center] at (8.4,-0.25)
 {Claim issued\\ 30-day limit\\ suspension bites};
 \node[anchor=north,text width=2.5cm,align=center] at (11.6,-0.25)
 {TCC judgment\\ 4 months--2 years};

 % artefact bands above
 \begin{scope}[on background layer]
 \fill[f21green!12,rounded corners=3pt] (-0.5,0.42) rectangle (3.9,1.5);
 \fill[f21amber!12,rounded corners=3pt] (4.3,0.42) rectangle (12.6,1.5);
 \end{scope}
 \node[anchor=west,text width=4.2cm,align=left,f21green] at (-0.45,0.96)
 {\textbf{Risk screen} (\S\ref{sec:risk-screen})\\ acts before a dispute exists};
 \node[anchor=west,text width=8.0cm,align=left,f21amber] at (4.4,0.96)
 {\textbf{Negotiation simulator} (\S\ref{sec:orchestration}) acts on a dispute\\ that has already been raised};
\end{tikzpicture}}
\caption[The procurement dispute lifecycle]{The procurement dispute lifecycle
under the Procurement Act 2023, and where this project's two artefacts sit on it.
The statutory windows are short; the litigation that follows is not. The horizontal
axis is ordinal, not to scale.}
\label{fig:lifecycle}
\end{figure}

This project was carried out with Fusion21, a social enterprise running
procurement frameworks for housing providers, local authorities and NHS bodies,
which sees the same disputes recur across its member base. Fusion21 supplied four
internal documents, referred to throughout as Docs 1--4: the interests and drivers
of each party, the escalation ladder from informal engagement to litigation, the
negotiation theory they apply, and three worked examples. Doc 4 makes a point that
quietly shapes the whole project: litigated cases are not typical, because most
procurement disputes settle before proceedings are issued. Every judgment used in
this evaluation is therefore a survivor of that filtering, and any claim built on
the resulting corpus inherits its bias.

\section{Subject area and related work}
\label{sec:related}

The technical territory here is multi-agent LLM simulation: several instances of
a language model, each assigned a role, a set of interests and a private
fallback position, made to act out a structured multi-party interaction. Three
strands of prior work are directly relevant.

\subsection*{Multi-agent LLM frameworks}

Once instruction-following models could handle few-shot task specification
\cite{brown2020language} and expose their intermediate reasoning on request
\cite{wei2022chain}, the prompt itself became a kind of programming surface. From
there it is a short step to running several copies of one model under different
role prompts and getting something that looks like a division of labour. That
is essentially what CAMEL \cite{li2023camel} does, pairing an assistant and a
user agent under an inception-prompting scheme, and what Generative Agents
\cite{park2023generative} extended by giving twenty-five simulated townspeople
memory, reflection and planning. AutoGen \cite{wu2023autogen} and MetaGPT
\cite{hong2024metagpt} took the same underlying idea and pointed it at
engineering workflows instead, and MetaGPT's authors argue specifically that
what stops multi-agent conversation from collapsing into noise is encoding
standard operating procedures directly into the prompts. That claim matters
here: the Court agent's verification protocol described in
Section~\ref{sec:court-design} is exactly this kind of encoded procedure, closer
in spirit to ReAct's interleaving of reasoning with a defined action space
\cite{yao2023react} than to open-ended chat, and Chapter~\ref{ch:evaluation}
shows plainly what happens once the case shapes stop matching what it was
written for. Orchestration in this project runs on LangGraph \cite{langgraph},
which represents the interaction as an explicit state graph rather than a
free-form conversation, a choice defended in Section~\ref{sec:orchestration}.

\subsection*{LLMs as negotiators}

Negotiation makes a good testbed because it comes with structure, an outcome, and
an existing literature. Fu et al.\ \cite{fu2023negotiation} set two models
bargaining over a balloon and let them improve through self-play with AI feedback.
Bianchi et al.'s NegotiationArena \cite{bianchi2024negotiation} built
resource-allocation, buy--sell and ultimatum games into a benchmark and found LLM
negotiators could be talked into bad deals by tactics like feigned desperation.
Closer to the setting here, Abdelnabi et al.\ \cite{abdelnabi2024negotiation}
built multi-party negotiations around conflicting private goals, scoring agents on
whether the deal reached satisfied their own constraints. What unites all three is
a self-contained environment: the payoff is defined by the game, so correctness is
decidable. Procurement disputes offer no such luxury. The governing framework sits
outside the simulation, is contested, and lives in statute and case law, so the
``correct'' outcome is whatever a court would have held, which is usually unknown.

The normative framework the agents reason with comes from Fisher and Ury's
\emph{Getting to Yes} \cite{fisher2011getting} (separate interests from positions,
know your BATNA) and Raiffa on asymmetric negotiation \cite{raiffa1982art}. Doc 3
adds a constraint from that literature which shapes the system's whole outcome
vocabulary: in award-stage disputes you cannot split the difference, because the
available remedies are fixed by statute.

\subsection*{Legal reasoning and the fabrication problem}

On legal benchmarks, language models do reasonably well. GPT-4 passed the Uniform
Bar Examination \cite{katz2024gpt4}, and LegalBench \cite{guha2023legalbench}
assembled 162 tasks showing real, if uneven, legal reasoning. The evidence that
matters more here cuts the other way. Dahl et al.\ \cite{dahl2024legal} profiled
legal hallucination and found error rates of 58--82\% on verifiable questions
about federal cases, with models rarely able to tell when they did not know.
Magesh et al.\ \cite{magesh2024hallucination} tested commercial
retrieval-augmented legal tools built specifically to avoid this and still
measured hallucination above 17\%. Ji et al.\ \cite{ji2023survey} put it plainly:
fluent, confident, wrong output is the default failure mode, not a rare edge case.

That is the difficulty in one sentence. A plausible-looking compliance finding on
a procurement dispute is not merely useless but actively harmful, precisely
because it is persuasive, unless something separately guarantees that the numbers
and provisions it cites came from the case in front of it. Retrieval
\cite{lewis2020rag} was the obvious candidate and was seriously considered;
Section~\ref{sec:alternatives} explains why it was scoped out and what was built
instead.

\section{The gap this project addresses}
\label{sec:gap}

Multi-agent negotiation research is typically scored against a payoff function
that the environment itself defines. Legal-LLM research is typically scored
against benchmark answers a human has already labelled. A simulated procurement
dispute has neither luxury: the artefact hands down a legal judgment on facts
that a real court either decided differently, decided on a different question, or
never decided. So the question this project organised itself around was not
\emph{can agents negotiate a procurement dispute}, which turned out to be the easy
part, but how you evaluate what they produce and how you stop them inventing the
evidence they reason from.

\section{Aim, research questions and objectives}
\label{sec:objectives}

\textbf{Aim.} To build and evaluate a multi-agent LLM system that simulates
procurement dispute resolution under the Procurement Act 2023 modelling the
interests, constraints and fallback positions of the contracting authority, the
aggrieved bidder and the court and, alongside it, to establish an evaluation
methodology suited to a domain where ground truth is scarce and fabrication is
the main risk.

The research questions carried over from the project plan, sharpened by
what the work actually turned up, are:

\begin{description}
 \item[RQ1] Without fine-tuning, can prompt engineering and structured context
 injection alone produce agents that reliably represent each party's distinct
 interests, legal constraints and BATNA?
 \item[RQ2] Which negotiation protocol lets autonomous agents reach legally
 grounded resolutions, and which single component of that protocol is
 actually doing the work?
 \item[RQ3] Can an LLM court agent be made to verify arithmetic when a case
 supplies one and refuse to invent one when it doesn't and does that
 constraint survive contact with dispute shapes it was never designed
 against?
 \item[RQ4] Given that most real disputes settle confidentially and never
 produce a ground-truth label, how should outcomes even be evaluated for
 reliability, fabrication and alignment with what actually happened?
\end{description}

These break down into six objectives (O1--O6), which give
Chapter~\ref{ch:methodology} its structure and are assessed against the evidence
in Table~\ref{tab:objectives}: specify the agents against Fusion21's own interest
taxonomies; build a deterministic orchestration protocol with schema-validated
output; design a court agent that verifies conditionally rather than by default;
build an ingestion path from real documents into structured scenarios; assemble a
verified corpus and evaluate against it with explicit baselines and ablations; and
deliver something Fusion21 can use, which turned out to mean a pre-award risk
screen rather than the simulator.

\section{Contributions}
\label{sec:contributions}

\begin{enumerate}
 \item \textbf{A conditional verification protocol for an LLM adjudicator}
 (Section~\ref{sec:court-design}). Exact arithmetic is only attempted when
 both a formula and every input it needs are present, and invented values are
 refused even when the agent tries to hedge them as illustrative. Four prompt
 revisions are documented, each written in response to a specific failure
 mode the previous version produced.
 \item \textbf{An evaluation methodology for a domain with almost no ground
 truth} (Chapter~\ref{ch:evaluation}). It combines agreement with real
 dispositions, three independent fabrication screens, instrumentation for
 structural compliance, a controlled prompt ablation with significance
 testing, and three baselines.
 \item \textbf{A negative result that limits how far the architecture claim can
 go} (Section~\ref{sec:baselines}): a naive zero-shot prompt matches the full
 pipeline on directional accuracy, while removing the court agent causes
 every run to deadlock. Whatever the architecture earns, it earns on
 resolution mechanics and process, not on accuracy.
 \item \textbf{Three distinct failure modes}, evidenced on named cases
 (Section~\ref{sec:fabrication}): numeric fabrication, narrative-premise
 fabrication, and legal-judgment divergence, together with the finding that
 hardening the adjudicator against the first does nothing about the
 second.
 \item \textbf{A verified corpus of 23 real UK procurement judgments}, rendered
 as structured scenarios with each disposition checked individually, released
 alongside the project.
 \item \textbf{A pre-award challenge risk screen} (Section~\ref{sec:risk-screen}),
 traced rule by rule back to Fusion21's own practical-considerations tables
 and the only part of this system with an observable ground truth in
 principle.
\end{enumerate}

\section{Report structure}

Chapter~\ref{ch:methodology} describes the system and defends the design choices
against the alternatives that were considered and rejected along the way.
Chapter~\ref{ch:evaluation} presents the evaluation, including the results that
do not flatter the artefact. Chapter~\ref{ch:conclusion} returns to the
objectives, reflects on how the project actually went, and sets out what should
be built next.